\documentclass[letterpaper,10pt]{article}

\usepackage{latexsym}
\usepackage[empty]{fullpage}
\usepackage{titlesec}
\usepackage{marvosym}
\usepackage[usenames,dvipsnames]{color}
\usepackage{verbatim}
\usepackage{enumitem}
\usepackage[hidelinks]{hyperref}
\usepackage{fancyhdr}
\usepackage[english]{babel}
\usepackage{tabularx}
\input{glyphtounicode}

\usepackage[sfdefault]{roboto}

\pagestyle{fancy}
\fancyhf{}
\fancyfoot{}
\renewcommand{\headrulewidth}{0pt}
\renewcommand{\footrulewidth}{0pt}
\setlength{\footskip}{5pt}

\addtolength{\oddsidemargin}{-0.5in}
\addtolength{\evensidemargin}{-0.5in}
\addtolength{\textwidth}{1in}
\addtolength{\topmargin}{-.5in}
\addtolength{\textheight}{1.0in}

\urlstyle{same}
\raggedbottom
\setlength{\tabcolsep}{0in}

% Text justification settings
\setlength{\parindent}{0pt}
\setlength{\parskip}{7pt}

\titleformat{\section}{\Large\bfseries\scshape\raggedright}{}{0em}{}[\titlerule]

\pdfgentounicode=1

\newcommand{\resumeItem}[1]{\item\small{#1}}
\newcommand{\resumeSubheading}[4]{
\vspace{-1pt}\item
  \begin{tabular*}{0.97\textwidth}[t]{l@{\extracolsep{\fill}}r}
    \textbf{#1} & #2 \\
    \textit{#3} & \textit{#4} \\
  \end{tabular*}\vspace{-7pt}
}
\renewcommand\labelitemii{$\vcenter{\hbox{\tiny$\bullet$}}$}
\newcommand{\resumeSubHeadingList}{\begin{itemize}[leftmargin=0.15in, label={}]}
\newcommand{\resumeSubHeadingListEnd}{\end{itemize}}

\begin{document}

\begin{center}
  \textbf{\Huge Gabriel Frigo} \\ \vspace{4pt}
  \small +55 11 95570-2000 $|$ \href{mailto:gabriel.frigo4@gmail.com}{gabriel.frigo4@gmail.com} $|$
  \href{https://linkedin.com/in/gabriel-frigo-b6727b275/}{linkedin.com/in/gabriel-frigo} $|$
  \href{https://github.com/GabrielFrigo4}{github.com/GabrielFrigo4} $|$
  \href{https://gabrielfrigo.dev.br}{gabrielfrigo.dev.br}
\end{center}

\vspace{6pt}

\textbf{Subject: Cover Letter \& Application -- Software Engineering}

Dear Engineering and Recruiting Team,

I am writing to express my interest in \textbf{Software Engineering, Systems, and Infrastructure} opportunities. I am a \textbf{Computer Science student at UFABC} whose profile bridges academic mathematical depth, pragmatic software engineering, low-level systems, and cloud infrastructure.

My technical mindset navigates seamlessly between \textbf{academic theory and real-world production}: I find beauty in the rigor of \textbf{Pure Mathematics and Operations Research} (applying Graph Theory, Network Flows, and Combinatorics) just as much as in modern applied engineering (IaC with OpenTofu, Go microservices, Docker/Podman, and cloud architecture). Likewise, I embrace standard enterprise tools without losing appreciation for hands-on systems craftsmanship---from writing custom shell scripts and configuring native \textbf{FreeBSD with Jails (Bastille)} environments to writing high-performance code in \textbf{C/C++, Lua, Got, and POSIX standards}.

On the practical side, I developed \textbf{OptiLaser}---a logistics decision and operations research engine (VRPTW) solving complex combinatorial optimization in $<3$s built in \textbf{Go and CGo} with \textbf{Google OR-Tools}, cross-platform (\textbf{FreeBSD/Linux}), featuring an embedded \textbf{PocketBase} database and IaC provisioning via \textbf{OpenTofu} on \textbf{Magalu Cloud} running on \textbf{Podman with Rocky Linux}. I also built a \textbf{concurrent HTTP server in C} from scratch over POSIX/BSD Sockets, and I am proficient in systems languages such as \textbf{C, C++, Rust, and Assembly}, which I apply in leading projects at the \textbf{Green Team Hacker Club} and developing software for competitive robotics.

My problem-solving skills have been forged through top-tier performance in competitions, such as the \textbf{Gold Medal (13th national rank) at OBMEP}, the \textbf{SBC Programming Marathon}, and the \textbf{ICPC}, alongside undergraduate research in \textbf{Network Flows (CNPq)} and \textbf{Ramsey Theory (PICME)}. Combining a strong \textbf{self-directed learning} drive with an \textbf{AI-First mindset}, I approach Artificial Intelligence not merely as a code generator, but as an intellectual partner for continuous learning via the \textbf{Socratic method}---questioning assumptions, rigorously analyzing architectural trade-offs, and accelerating my technical mastery.

I believe this versatility---uniting mathematical modeling, low-level systems, scalable infrastructure, and rapid learning---qualifies me to make an immediate impact on high-impact technical challenges. I am available for an interview at your convenience.

Sincerely,

\textbf{Gabriel Frigo}

\end{document}
